\documentclass[10pt,a4paper]{amsart}
\usepackage[margin=19mm]{geometry}
\usepackage{amsmath,amssymb,mathtools}
\usepackage[scr=rsfs,cal=euler]{mathalfa}
\usepackage{xfrac}
\usepackage[unicode,hidelinks,bookmarks=false]{hyperref}
\newcommand{\ie}{i.e.\ }
\newcommand{\cf}{cf.\ }
\newcommand{\resp}{respectively\ }
\newcommand{\Subsection}{\subsection}
\providecommand{\nobreakdash}{\nobreak\hbox{-}}
\setlength{\emergencystretch}{2em}
\multlinegap0pt
\usepackage{amssymb}
\usepackage{bm}
\usepackage{mathtools}
\usepackage{algorithm}\usepackage{algpseudocode}
\usepackage{tikz}
\usepackage{xcolor}
\newtheorem{lemma}{Lemma}
\newcommand{\bx}{\boldsymbol{x}}
\title{A Reciprocity-Law-Compliant Photoacoustic Forward-Adjoint Operator}
\author{Ashkan Javaherian}
\date{Source: arXiv:2604.04022v3 (CC BY 4.0)}
\begin{document}
\maketitle

\noindent\emph{Source and licence.} A Reciprocity-Law-Compliant Photoacoustic Forward-Adjoint Operator. Version arXiv:2604.04022v3 (CC BY 4.0), distributed by arXiv under the Creative Commons Attribution 4.0 International licence, http://creativecommons.org/licenses/by/4.0/, as recorded in the arXiv licence metadata for this exact version and shown on the abstract page https://arxiv.org/abs/2604.04022v3. Full licence notice: \texttt{licenses/arxiv-2604.04022-CC-BY-4.0.txt}.\par\smallskip

\noindent\textit{Editorial scope.} Counted unit: the article's lemma eq:adjoint-photoacoustic (source0.tex lines 335--347). The article's complete original proof of it is delivered with it (source0.tex lines 349--371, one \texttt{proof} environment of 23 lines). No mathematics is written, repaired or re-derived: the statement and the proof are the article's own lines, in the article's own notation, and no retained line is substituted or re-flowed.  Retained uncounted as the source-local context the proof needs: lines 277--282.  Nothing was omitted from any retained span, so the package carries no omission record.  No retained line contains a citation, so the wrapper carries no bibliography and no bibliography entry.  No retained line contains a figure environment, an image-inclusion command or drawing code, so no image is delivered and the package declares no asset dependency.  Editorial formatting only, in addition: an AMS \texttt{amsart} preamble that restates the article's own declarations and macros used by the retained lines, namely the environments \texttt{lemma} on separate counters, together with the standard packages those lines need.  The retained environments print in this document as Lemma 1 in order of appearance, whereas the article prints the same items with the numbering of its own full document.  The complete native source of the article as distributed in the arXiv e-print for this exact version is bundled unchanged as \texttt{sources/197855/source0.tex}.
\begin{align} \label{eq:photoacoustic-op}
    \begin{split}
        \mathcal{Q}: C_0^{\infty}(\Omega) &\rightarrow C_0^{\infty}(\mathbb{R}^d \times \mathbb{R}), \\
        \mathcal{Q}[p_0] &= s_{\mathcal{Q}},
    \end{split}
\end{align}

\begin{lemma}
 The action of the adjoint of the operator \(\mathcal{Q} \), defined by Eq.~\eqref{eq:photoacoustic-op}, on any test function 
    \( f^{s_{\mathcal{Q}}} \in C_0^\infty(\mathbb{R}^d \times \mathbb{R}) \) is given by
    \begin{align}  \label{eq:adjoint-photoacoustic}
        \begin{split}
        &\mathcal{Q}^*: C_0^\infty(\mathbb{R}^d \times \mathbb{R}) 
        \rightarrow C_0^\infty(\Omega), \\
        &\mathcal{Q}^*[f^{s_{\mathcal{Q}}}]
        = - \frac{1}{c(\bx)^2}\,  \,   \delta_q(t) \, \frac{\partial  f^{s_{\mathcal{Q}} } }{\partial t}
       (\bx,t) \rvert_{t=0}.
       \end{split}
    \end{align}
\end{lemma}

\begin{proof}
The adjoint operator \( \mathcal{Q}\) and its adjoint, with respect to the standard 
\(L^2\) bilinear form in  \(C_0^\infty(\Omega)\) and \(C_0^\infty(\mathbb{R}^d \times \mathbb{R})\), should satisfy
\begin{align} \label{eq:photoacoustic-adj-1}
        \int_{\mathbb{R}^d} d\bx \, \int_{\mathbb{R}} dt \,
        f^{s_{\mathcal{Q}}}(\bx,t)  \,
        \frac{1}{c(\bx)^2} f_0(\bx)  \,  \frac{d}{dt}  \delta_q(t)
        = \int_{\Omega} d\bx \, \int_{\mathbb{R}} dt \, \mathcal{Q}^*[f^{s_{\mathcal{Q}}}]  \, f_0(\bx)
\end{align}
for any  \( f_0 \in C_0^\infty(\Omega)\)~and~\(f^{s_{\mathcal{Q}}} \in C_0^\infty(\mathbb{R}^d \times \mathbb{R})\).

Applying integration by parts in time to the left-hand side of 
\eqref{eq:photoacoustic-adj-1} yields
\begin{align} \label{eq:photoacoustic-adj-2}
    \int_{\Omega} \, d\bx \,   \int_{\mathbb{R}}   dt  \,  \Big[ - \frac{1}{ c(\bx)^2} 
    \delta_q(t) \
    \frac{\partial f^{s_{\mathcal{Q}}} }{\partial t} (\bx,t)  \Big]
    f(\bx),
\end{align}
where we have also restricted the domain of spatial integration to the support of \(f\).

Comparing the integral expression in \eqref{eq:photoacoustic-adj-2} with the right-hand side of Eq.~\eqref{eq:photoacoustic-adj-1} immediately identifies the bracketed term in \eqref{eq:photoacoustic-adj-2} as \(\mathcal{Q}^*[f^{s_{\mathcal{Q}}}]\), which proves the expression in \eqref{eq:adjoint-photoacoustic}.
\end{proof}

\end{document}
